\documentclass{article}%
\usepackage[T1]{fontenc}%
\usepackage[utf8]{inputenc}%
\usepackage{lmodern}%
\usepackage{textcomp}%
\usepackage{lastpage}%
\usepackage{geometry}%
\geometry{margin=1in,headheight=14pt,headsep=25pt}%
\usepackage{hyperref}%
\usepackage{enumitem}%
\usepackage{fancyhdr}%
\usepackage{xcolor}%
%

            \hypersetup{
                colorlinks=true,
                linkcolor=blue,
                filecolor=magenta,
                urlcolor=blue,
            }
            
            \pagestyle{fancy}
            \fancyhf{}
            \rhead{Generated on \today}
            \lhead{Course Summary}
            \cfoot{\thepage}
            
            % Configure itemize settings globally
            \setlist[itemize]{nosep}  % Reduce vertical spacing between items
            \setlist[itemize]{leftmargin=*}  % Align with left margin
        %
\title{Linear Programming Course Summary}%
\author{Generated by Scribe.AI}%
\date{\today}%
%
\begin{document}%
\normalsize%
\maketitle%
\section{Key Terms}%
\label{sec:KeyTerms}%
\subsection{Linear Programming Fundamentals}%
\label{subsec:LinearProgrammingFundamentals}%
\begin{itemize}[leftmargin=*]%
\item \href{https://www.math.purdue.edu/~yipn/421/2024-08-27-ExSimplex.pdf#page=1}{\textbf{Slack Variable}}: A variable added to a less-than-or-equal-to inequality constraint to transform it into an equality constraint. It represents the difference between the left-hand side and right-hand side of the inequality.%
\item \href{https://www.math.purdue.edu/~yipn/421/2024-08-27-ExSimplex.pdf#page=4}{\textbf{Feasible Region}}: The set of all points that satisfy all the constraints of a linear programming problem.%
\item \href{https://www.math.purdue.edu/~yipn/421/2024-08-27-ExSimplex.pdf#page=5}{\textbf{Vertex}}: A point where two or more constraint boundaries intersect in the feasible region of a linear programming problem. In the Simplex method, optimal solutions are found at vertices.%
\item \href{https://www.math.purdue.edu/~yipn/421/2024-08-27-ExSimplex.pdf#page=6}{\textbf{Basic Variable}}: In the Simplex method, a variable that is expressed in terms of non-basic variables in a dictionary. In a basic feasible solution, basic variables are non-negative.%
\item \href{https://www.math.purdue.edu/~yipn/421/2024-08-27-ExSimplex.pdf#page=6}{\textbf{Non-basic Variable}}: In the Simplex method, a variable set to zero in a given iteration of the algorithm. The values of the basic variables are determined by the values of the non-basic variables.%
\item \href{https://www.math.purdue.edu/~yipn/421/2024-08-27-ExSimplex.pdf#page=9}{\textbf{Pivot Operation}}: The process of exchanging a basic and a non-basic variable in the Simplex method, leading to a new basic feasible solution.%
\item \href{https://www.math.purdue.edu/~yipn/421/2024-08-27-ExSimplex.pdf#page=26}{\textbf{Dictionary}}: A representation of a linear programming problem where basic variables are expressed as functions of non-basic variables. It is used in the Simplex method to systematically move from one vertex to another.%
\item \href{https://www.math.purdue.edu/~yipn/421/2024-08-27-ExSimplex.pdf#page=32}{\textbf{Auxiliary Problem}}: A modified linear programming problem introduced to find an initial feasible solution when the origin is not feasible in the original problem. It involves adding a new variable to shift the feasible region.%
\item \href{https://www.math.purdue.edu/~yipn/421/2024-09-05-SimplexEfficiency.pdf#page=4}{\textbf{Largest-Coefficient Rule}}: A pivot rule in the Simplex method that selects the entering variable with the largest coefficient in the objective function row of the dictionary.%
\item \href{https://www.math.purdue.edu/~yipn/421/2024-09-05-SimplexEfficiency.pdf#page=6}{\textbf{Largest-Increase Rule}}: A pivot rule in the Simplex method that selects the entering variable that results in the largest increase in the objective function value.%
\item \href{https://www.math.purdue.edu/~yipn/421/2024-09-05-SimplexEfficiency.pdf#page=2}{\textbf{Klee-Minty Example}}: A worst-case example for the Simplex method, demonstrating that the number of iterations can grow exponentially with the problem size.%
\item \href{https://www.math.purdue.edu/~yipn/421/2024-09-05-SimplexEfficiency.pdf#page=10}{\textbf{Bland's Rule}}: A pivot rule in the Simplex method designed to prevent cycling, ensuring termination.%
\item \href{https://www.math.purdue.edu/~yipn/421/2024-09-05-SimplexEfficiency.pdf#page=10}{\textbf{Lexicographic Method}}: A method for choosing the leaving variable in the Simplex method that ensures termination by preventing cycling.%
\item \href{https://www.math.purdue.edu/~yipn/421/2024-09-05-SimplexEfficiency.pdf#page=10}{\textbf{Randomized Pivot Rule}}: A pivot rule that randomly permutes the order of variables before applying another rule, such as Bland's rule.%
\item \href{https://www.math.purdue.edu/~yipn/421/2024-09-05-SimplexEfficiency.pdf#page=11}{\textbf{Smoothed Complexity}}: A framework for analyzing the complexity of algorithms by considering the average-case performance after small random perturbations to the input.%
\item \href{https://www.math.purdue.edu/~yipn/421/2024-09-12-Duality.pdf#page=1}{\textbf{Primal Problem}}: A linear programming problem that aims to maximize an objective function subject to a set of constraints and non-negativity conditions.%
\item \href{https://www.math.purdue.edu/~yipn/421/2024-10-03-DualSimplexMatrix.pdf#page=9}{\textbf{Entering Variable}}: In the simplex method, the non-basic variable selected to become a basic variable in the next iteration. The selection is typically based on improving the objective function value.%
\item \href{https://www.math.purdue.edu/~yipn/421/2024-10-03-DualSimplexMatrix.pdf#page=10}{\textbf{Leaving Variable}}: In the simplex method, the basic variable selected to become a non-basic variable in the next iteration. The selection ensures that the resulting solution remains feasible.%
\item \href{https://www.math.purdue.edu/~yipn/421/2024-10-03-DualSimplexMatrix.pdf#page=15}{\textbf{Elementary Matrix}}: A matrix obtained by performing a single elementary row operation (such as swapping two rows, multiplying a row by a non-zero scalar, or adding a multiple of one row to another) on an identity matrix. Used in updating the basis matrix in the simplex method.%
\item \href{https://www.math.purdue.edu/~yipn/421/2024-10-28-Farkas.pdf#page=3}{\textbf{Fredholm Alternatives}}: A statement of the two mutually exclusive possibilities regarding the solvability of the system $Ax = b$.%
\item \href{https://www.math.purdue.edu/~yipn/421/2024-12-03-Duality.pdf#page=3}{\textbf{Graphical Method}}: A method for solving linear programming problems with two variables by visualizing the feasible region and objective function on a graph. The optimal solution is found at a vertex of the feasible region.%
\end{itemize}

%
\subsection{Duality Theory}%
\label{subsec:DualityTheory}%
\begin{itemize}[leftmargin=*]%
\item \href{https://www.math.purdue.edu/~yipn/421/2024-09-12-Duality.pdf#page=1}{\textbf{Dual Problem}}: A linear programming problem derived from the primal problem, aiming to minimize an objective function subject to a different set of constraints and non-negativity conditions. The optimal solution of the dual problem provides valuable information about the primal problem.%
\item \href{https://www.math.purdue.edu/~yipn/421/2024-09-12-Duality.pdf#page=2}{\textbf{Weak Duality}}: A theorem stating that the objective function value of any feasible solution to the primal problem is always less than or equal to the objective function value of any feasible solution to the dual problem.%
\item \href{https://www.math.purdue.edu/~yipn/421/2024-09-12-Duality.pdf#page=6}{\textbf{Negative Transpose Property}}: The property that the coefficient matrix of the dual problem is the negative transpose of the coefficient matrix of the primal problem. This property is preserved throughout the simplex method iterations.%
\item \href{https://www.math.purdue.edu/~yipn/421/2024-09-12-Duality.pdf#page=11}{\textbf{Strong Duality Theorem}}: A theorem stating that if a linear programming problem has an optimal solution, then its dual problem also has an optimal solution, and the optimal objective function values of both problems are equal.%
\item \href{https://www.math.purdue.edu/~yipn/421/2024-09-12-Duality.pdf#page=5}{\textbf{Complementary Slackness}}: A condition stating that for optimal primal and dual solutions, the product of each primal variable and its corresponding dual slack variable is zero, and the product of each dual variable and its corresponding primal slack variable is zero.%
\item \href{https://www.math.purdue.edu/~yipn/421/2024-09-17-DualityGeneral.pdf#page=2}{\textbf{Lagrange Multiplier}}: A scalar value associated with each constraint in a constrained optimization problem, used to incorporate the constraints into the objective function via the Lagrangian.%
\item \href{https://www.math.purdue.edu/~yipn/421/2024-09-17-DualityGeneral.pdf#page=2}{\textbf{Lagrangian Function}}: A function that combines the objective function and constraints of a constrained optimization problem using Lagrange multipliers. The Lagrangian is unconstrained in the original variables.%
\item \href{https://www.math.purdue.edu/~yipn/421/2024-09-17-DualityGeneral.pdf#page=5}{\textbf{Weak Duality Theorem}}: A theorem stating that the optimal value of the primal problem is always less than or equal to the optimal value of the dual problem.%
\item \href{https://www.math.purdue.edu/~yipn/421/2024-10-01-SensitivityDuality.pdf#page=10}{\textbf{Shadow Price}}: The rate of change in the optimal objective function value for a unit change in the right-hand side of a constraint.%
\item \href{https://www.math.purdue.edu/~yipn/421/2024-10-01-SensitivityDuality.pdf#page=11}{\textbf{Reduced Costs}}: The vector $z_N = (B^\textbackslash{}{-1\textbackslash{}}N)^T c_B - c_N$, representing the change in the objective function value per unit increase in a non-basic variable. A non-negative reduced cost vector indicates optimality.%
\item \href{https://www.math.purdue.edu/~yipn/421/2024-10-01-SensitivityDuality.pdf#page=3}{\textbf{Complementary Slackness Conditions}}: Conditions that relate the primal and dual optimal solutions. Specifically, for each constraint, either the constraint is binding (equality holds) or the corresponding dual variable is zero (and vice versa).%
\item \href{https://www.math.purdue.edu/~yipn/421/2024-10-01-SensitivityDuality.pdf#page=1}{\textbf{Sensitivity Analysis}}: The study of how changes in the parameters of a linear programming problem (e.g., objective function coefficients, constraint coefficients, or right-hand side values) affect the optimal solution.%
\item \href{https://www.math.purdue.edu/~yipn/421/2024-10-01-SensitivityDuality.pdf#page=4}{\textbf{Basis Matrix ($B$)}}: A submatrix of the constraint matrix $A$ formed by the columns corresponding to the basic variables. Its inverse is used in the simplex method to update the solution.%
\item \href{https://www.math.purdue.edu/~yipn/421/2024-10-01-SensitivityDuality.pdf#page=4}{\textbf{Non-basis Matrix ($N$)}}: A submatrix of the constraint matrix $A$ formed by the columns corresponding to the non-basic variables.%
\item \href{https://www.math.purdue.edu/~yipn/421/2024-10-01-SensitivityDuality.pdf#page=5}{\textbf{Simplex Tableau}}: A tabular representation of the linear programming problem used in the simplex method to track the values of the variables and the objective function during iterations.%
\item \href{https://www.math.purdue.edu/~yipn/421/2024-10-22-Convexity.pdf#page=12}{\textbf{Farkas' Lemma}}: The system $Ax \textbackslash{}le b$ has no solution if and only if there exists a vector $y \textbackslash{}ge 0$ such that $y^*T A = 0$ and $y^*T b < 0$.%
\end{itemize}

%
\subsection{Advanced Linear Programming Techniques}%
\label{subsec:AdvancedLinearProgrammingTechniques}%
\begin{itemize}[leftmargin=*]%
\item \href{https://www.math.purdue.edu/~yipn/421/2024-09-05-SimplexEfficiency.pdf#page=12}{\textbf{Interior Method}}: A class of algorithms for linear programming that traverse the interior of the feasible region, in contrast to the Simplex method which follows the boundary.%
\item \href{https://www.math.purdue.edu/~yipn/421/2024-10-03-DualSimplexMatrix.pdf#page=2}{\textbf{Dual Simplex Method}}: A variant of the simplex method that starts with a dual feasible solution and iteratively moves towards primal feasibility while maintaining dual feasibility, ultimately finding an optimal solution for both the primal and dual problems.%
\item \href{https://www.math.purdue.edu/~yipn/421/2024-10-03-DualSimplexMatrix.pdf#page=5}{\textbf{Dual Based Phase I Algorithm}}: A method used to find an initial feasible solution for the dual problem in linear programming, particularly when the origin is not dual feasible. It involves modifying the dual problem to ensure feasibility at the origin, then applying the dual simplex method.%
\end{itemize}

%
\subsection{Linear Programming Applications}%
\label{subsec:LinearProgrammingApplications}%
\begin{itemize}[leftmargin=*]%
\item \href{https://www.math.purdue.edu/~yipn/421/2024-09-19-DualityExamples.pdf#page=2}{\textbf{Diet Problem}}: A linear programming problem where the objective is to minimize the cost of a diet while meeting minimum nutritional requirements.%
\item \href{https://www.math.purdue.edu/~yipn/421/2024-09-19-DualityExamples.pdf#page=2}{\textbf{Nutrient}}: A substance that provides nourishment essential for growth and the maintenance of life.%
\item \href{https://www.math.purdue.edu/~yipn/421/2024-09-19-DualityExamples.pdf#page=2}{\textbf{Food}}: An edible substance that provides nutrients.%
\item \href{https://www.math.purdue.edu/~yipn/421/2024-09-19-DualityExamples.pdf#page=1}{\textbf{Constraint Matrix}}: A matrix representing the coefficients of the variables in the constraints of a linear programming problem.%
\item \href{https://www.math.purdue.edu/~yipn/421/2024-09-19-DualityExamples.pdf#page=1}{\textbf{Decision Variables}}: Variables representing the quantities to be determined in an optimization problem.%
\item \href{https://www.math.purdue.edu/~yipn/421/2024-09-19-DualityExamples.pdf#page=1}{\textbf{Objective Function}}: The function to be minimized or maximized in an optimization problem.%
\item \href{https://www.math.purdue.edu/~yipn/421/2024-10-28-ApplsLinProg.pdf#page=1}{\textbf{Monthly Demands}}: The predicted sales of ice cream for each month of the year.%
\item \href{https://www.math.purdue.edu/~yipn/421/2024-10-28-ApplsLinProg.pdf#page=2}{\textbf{Average Demand}}: The average monthly sales of ice cream over the entire year.%
\item \href{https://www.math.purdue.edu/~yipn/421/2024-10-28-ApplsLinProg.pdf#page=3}{\textbf{Cost of Changing Production}}: The cost incurred by the ice cream manufacturer for altering its production levels from one month to the next.%
\item \href{https://www.math.purdue.edu/~yipn/421/2024-10-28-ApplsLinProg.pdf#page=3}{\textbf{Cost of Storage}}: The cost incurred by the ice cream manufacturer for storing surplus ice cream from one month to the next.%
\item \href{https://www.math.purdue.edu/~yipn/421/2024-10-28-ApplsLinProg.pdf#page=3}{\textbf{Production Level}}: The amount of ice cream produced in a given month.%
\item \href{https://www.math.purdue.edu/~yipn/421/2024-10-28-ApplsLinProg.pdf#page=3}{\textbf{Surplus}}: The excess amount of ice cream remaining at the end of a month after meeting the demand.%
\item \href{https://www.math.purdue.edu/~yipn/421/2024-10-28-ApplsLinProg.pdf#page=9}{\textbf{Markowitz Portfolio Selection Model}}: A financial model that aims to optimize investment portfolios by maximizing expected return while minimizing risk.%
\item \href{https://www.math.purdue.edu/~yipn/421/2024-10-28-ApplsLinProg.pdf#page=9}{\textbf{Annual Return}}: The return on investment for a given asset over a year.%
\item \href{https://www.math.purdue.edu/~yipn/421/2024-10-28-ApplsLinProg.pdf#page=9}{\textbf{Proportion of Investment}}: The fraction of a portfolio allocated to a specific asset.%
\item \href{https://www.math.purdue.edu/~yipn/421/2024-10-28-ApplsLinProg.pdf#page=10}{\textbf{Fluctuation}}: The deviation of an asset's return from its expected return.%
\item \href{https://www.math.purdue.edu/~yipn/421/2024-10-28-ApplsLinProg.pdf#page=13}{\textbf{Portfolio Variance}}: A measure of the risk associated with a portfolio, calculated as the weighted sum of the covariances between all pairs of assets.%
\item \href{https://www.math.purdue.edu/~yipn/421/2024-10-28-ApplsLinProg.pdf#page=14}{\textbf{Covariance}}: A measure of the joint variability of two random variables.%
\item \href{https://www.math.purdue.edu/~yipn/421/2024-10-28-ApplsLinProg.pdf#page=17}{\textbf{Smallest Enclosing Ball Problem}}: The problem of finding the smallest hypersphere that encloses a given set of points in a $d$-dimensional space.%
\item \href{https://www.math.purdue.edu/~yipn/421/2024-10-28-ApplsLinProg.pdf#page=19}{\textbf{Maximum Weight Matching}}: A combinatorial optimization problem that involves finding a matching in a weighted graph such that the sum of the weights of the edges in the matching is maximized.%
\item \href{https://www.math.purdue.edu/~yipn/421/2024-10-28-ApplsLinProg.pdf#page=22}{\textbf{Machine Scheduling}}: A combinatorial optimization problem that involves assigning jobs to machines to minimize the total completion time.%
\item \href{https://www.math.purdue.edu/~yipn/421/2024-10-28-ApplsLinProg.pdf#page=24}{\textbf{Knapsack Problem}}: A combinatorial optimization problem that involves selecting a subset of items with weights and values to maximize the total value while not exceeding a given weight capacity.%
\item \href{https://www.math.purdue.edu/~yipn/421/2024-10-28-ApplsLinProg.pdf#page=25}{\textbf{Cutting Stock Problem}}: An optimization problem that involves determining how to cut large rolls of material into smaller rolls of various widths to minimize waste.%
\item \href{https://www.math.purdue.edu/~yipn/421/2024-10-28-ApplsLinProg.pdf#page=30}{\textbf{Sparse Solution}}: A solution to a system of linear equations that has a limited number of non-zero elements.%
\item \href{https://www.math.purdue.edu/~yipn/421/2024-10-28-ApplsLinProg.pdf#page=34}{\textbf{$l_1$ Norm}}: The sum of the absolute values of the elements of a vector.%
\end{itemize}

%
\subsection{Convex Optimization}%
\label{subsec:ConvexOptimization}%
\begin{itemize}[leftmargin=*]%
\item \href{https://www.math.purdue.edu/~yipn/421/2024-10-22-Convexity.pdf#page=1}{\textbf{Convex Set}}: A set $S \textbackslash{}subseteq \textbackslash{}mathbb\textbackslash{}{R\textbackslash{}}^m$ such that for any $z_1, z_2 \textbackslash{}in S$, $\textbackslash{}lambda z_1 + (1-\textbackslash{}lambda) z_2 \textbackslash{}in S$ for all $0 \textbackslash{}le \textbackslash{}lambda \textbackslash{}le 1$.%
\item \href{https://www.math.purdue.edu/~yipn/421/2024-10-22-Convexity.pdf#page=1}{\textbf{Convex Combination}}: A weighted sum of points $z_1, z_2, \textbackslash{}dots, z_n \textbackslash{}in \textbackslash{}mathbb\textbackslash{}{R\textbackslash{}}^m$ of the form $\textbackslash{}sum_\textbackslash{}{i=1\textbackslash{}}^n \textbackslash{}lambda_i z_i$, where $\textbackslash{}lambda_i \textbackslash{}ge 0$ for all $i$ and $\textbackslash{}sum_\textbackslash{}{i=1\textbackslash{}}^n \textbackslash{}lambda_i = 1$.%
\item \href{https://www.math.purdue.edu/~yipn/421/2024-10-22-Convexity.pdf#page=5}{\textbf{Convex Hull}}: The smallest convex set containing a given set $S \textbackslash{}subseteq \textbackslash{}mathbb\textbackslash{}{R\textbackslash{}}^m$, denoted as $\textbackslash{}text\textbackslash{}{Conv\textbackslash{}}(S) = \textbackslash{}bigcap \textbackslash{}\textbackslash{}{C \textbackslash{}text\textbackslash{}{ convex \textbackslash{}} | S \textbackslash{}subseteq C\textbackslash{}\textbackslash{}}$.%
\item \href{https://www.math.purdue.edu/~yipn/421/2024-10-22-Convexity.pdf#page=8}{\textbf{Caratheodory Theorem}}: If $z \textbackslash{}in \textbackslash{}text\textbackslash{}{Conv\textbackslash{}}(S)$, then $z$ can be written as a convex combination of at most $m+1$ points from $S$.%
\item \href{https://www.math.purdue.edu/~yipn/421/2024-10-22-Convexity.pdf#page=10}{\textbf{Separation Theorem}}: Let $P$ and $\textbackslash{}bar\textbackslash{}{P\textbackslash{}}$ be two disjoint nonempty polyhedra in $\textbackslash{}mathbb\textbackslash{}{R\textbackslash{}}^n$. Then there exist disjoint half-spaces $H$ and $\textbackslash{}bar\textbackslash{}{H\textbackslash{}}$ such that $P \textbackslash{}subset H$ and $\textbackslash{}bar\textbackslash{}{P\textbackslash{}} \textbackslash{}subset \textbackslash{}bar\textbackslash{}{H\textbackslash{}}$.%
\item \href{https://www.math.purdue.edu/~yipn/421/2024-10-22-Regression.pdf#page=15}{\textbf{Maximum Inscribed Circle in a Convex Polyhedron}}: Finding the largest circle that can fit entirely inside a convex polyhedron.%
\item \href{https://www.math.purdue.edu/~yipn/421/2024-10-22-Regression.pdf#page=22}{\textbf{Smallest Ball Problem}}: Finding the smallest radius ball that contains all given points in $\textbackslash{}mathbb\textbackslash{}{R\textbackslash{}}^d$.%
\end{itemize}

%
\subsection{Network Optimization}%
\label{subsec:NetworkOptimization}%
\begin{itemize}[leftmargin=*]%
\item \href{https://www.math.purdue.edu/~yipn/421/2024-11-05-NetFlows.pdf#page=1}{\textbf{Network Flow}}: A mathematical model representing the flow of commodities through a network, often formulated as a linear program.%
\item \href{https://www.math.purdue.edu/~yipn/421/2024-11-05-NetFlows.pdf#page=1}{\textbf{Node}}: A point in a network representing a source, sink, or intermediate point in the flow.%
\item \href{https://www.math.purdue.edu/~yipn/421/2024-11-05-NetFlows.pdf#page=1}{\textbf{Arc}}: A directed edge in a network connecting two nodes, representing a flow path with a capacity or cost.%
\item \href{https://www.math.purdue.edu/~yipn/421/2024-11-05-NetFlows.pdf#page=2}{\textbf{Supply}}: The amount of commodity available at a node in a network flow problem.%
\item \href{https://www.math.purdue.edu/~yipn/421/2024-11-05-NetFlows.pdf#page=2}{\textbf{Demand}}: The amount of commodity required at a node in a network flow problem.%
\item \href{https://www.math.purdue.edu/~yipn/421/2024-11-05-NetFlows.pdf#page=4}{\textbf{Incidence Matrix}}: A matrix representing the connections between nodes and arcs in a network, used in the linear programming formulation of network flow problems.%
\item \href{https://www.math.purdue.edu/~yipn/421/2024-11-05-NetFlows.pdf#page=6}{\textbf{Root Node}}: A designated node in a network, often used as a reference point for simplifying calculations in network flow problems.%
\item \href{https://www.math.purdue.edu/~yipn/421/2024-11-05-NetFlows.pdf#page=10}{\textbf{Spanning Tree}}: A tree that connects all nodes in a network, used in solving network flow problems.%
\item \href{https://www.math.purdue.edu/~yipn/421/2024-11-05-NetFlows.pdf#page=11}{\textbf{Tree Solution}}: A solution to a network flow problem where the flow along arcs in a spanning tree is zero.%
\item \href{https://www.math.purdue.edu/~yipn/421/2024-11-05-NetFlows.pdf#page=22}{\textbf{Reduced Cost}}: The cost of increasing the flow along a non-basic arc in a network flow problem.%
\item \href{https://www.math.purdue.edu/~yipn/421/2024-11-05-NetFlows.pdf#page=22}{\textbf{Primal Simplex Method}}: An algorithm used to solve linear programming problems, applied here to solve network flow problems by iteratively improving a tree solution.%
\item \href{https://www.math.purdue.edu/~yipn/421/2024-11-07-NetSimplex.pdf#page=1}{\textbf{Network}}: A graph consisting of a set of nodes (N) and a set of arcs (A).%
\item \href{https://www.math.purdue.edu/~yipn/421/2024-11-07-NetSimplex.pdf#page=2}{\textbf{Primal Flow}}: A flow assignment in a network that satisfies the constraints of the primal problem.%
\item \href{https://www.math.purdue.edu/~yipn/421/2024-11-07-NetSimplex.pdf#page=3}{\textbf{Dual Flow}}: A flow assignment in a network that satisfies the constraints of the dual problem.%
\item \href{https://www.math.purdue.edu/~yipn/421/2024-11-07-NetSimplex.pdf#page=3}{\textbf{Dual Slack}}: The difference between the cost of an arc and the difference in dual variables between its endpoints.%
\item \href{https://www.math.purdue.edu/~yipn/421/2024-11-07-NetSimplex.pdf#page=45}{\textbf{Two-Phased Network Simplex Method}}: A method that first finds a primal feasible solution using artificial costs and then uses the primal network simplex method to find an optimal solution.%
\item \href{https://www.math.purdue.edu/~yipn/421/2024-11-07-NetSimplex.pdf#page=45}{\textbf{Parametric Self-Dual Method}}: A method that intermingles primal and dual pivots to reduce a perturbation parameter from infinity to zero.%
\item \href{https://www.math.purdue.edu/~yipn/421/2024-11-14-NetAppls.pdf#page=1}{\textbf{Transportation Problem}}: A network flow problem where the objective is to minimize the cost of transporting goods from multiple sources to multiple destinations, subject to supply and demand constraints that may be inequalities.%
\item \href{https://www.math.purdue.edu/~yipn/421/2024-11-14-NetAppls.pdf#page=3}{\textbf{Dummy Node}}: A node added to a network flow model to handle inequality constraints in supply and demand, allowing for excess supply or unmet demand.%
\item \href{https://www.math.purdue.edu/~yipn/421/2024-11-14-NetAppls.pdf#page=9}{\textbf{Upper Bounded Transshipment Problem}}: A network flow problem where the flow along each arc is bounded above by a specified upper bound, in addition to supply and demand constraints.%
\item \href{https://www.math.purdue.edu/~yipn/421/2024-11-14-NetAppls.pdf#page=13}{\textbf{Shortest Path Problem}}: A graph problem where the objective is to find the path of minimum cost or distance between a source node and a destination node in a weighted graph.%
\item \href{https://www.math.purdue.edu/~yipn/421/2024-11-14-NetAppls.pdf#page=16}{\textbf{Bellman's Equation}}: A recursive equation used to solve the shortest path problem by iteratively computing the shortest distance from each node to a destination node.%
\item \href{https://www.math.purdue.edu/~yipn/421/2024-11-14-NetAppls.pdf#page=20}{\textbf{Dijkstra's Algorithm}}: A greedy algorithm for finding the shortest paths from a single source node to all other nodes in a graph with non-negative edge weights.%
\item \href{https://www.math.purdue.edu/~yipn/421/2024-11-19-MaxFlowMinCut.pdf#page=1}{\textbf{Max Flow}}: The maximum amount of flow that can be sent from a source node to a sink node in a network, subject to capacity constraints on the edges.%
\item \href{https://www.math.purdue.edu/~yipn/421/2024-11-19-MaxFlowMinCut.pdf#page=2}{\textbf{Cut}}: A partition of the nodes in a network into two disjoint sets, one containing the source node and the other containing the sink node.%
\item \href{https://www.math.purdue.edu/~yipn/421/2024-11-19-MaxFlowMinCut.pdf#page=2}{\textbf{Cut Capacity}}: The sum of the capacities of all edges that cross a cut from the source set to the sink set.%
\item \href{https://www.math.purdue.edu/~yipn/421/2024-11-19-MaxFlowMinCut.pdf#page=4}{\textbf{Max-Flow Min-Cut Theorem}}: The theorem stating that the maximum flow in a network is equal to the minimum capacity of any cut in the network.%
\item \href{https://www.math.purdue.edu/~yipn/421/2024-11-19-MaxFlowMinCut.pdf#page=10}{\textbf{Augmenting Path}}: A path in a network from the source to the sink along which flow can be increased.%
\item \href{https://www.math.purdue.edu/~yipn/421/2024-11-19-MaxFlowMinCut.pdf#page=10}{\textbf{Ford-Fulkerson Algorithm}}: An algorithm for finding the maximum flow in a network by iteratively finding and augmenting along augmenting paths.%
\end{itemize}

%
\subsection{Integer Programming}%
\label{subsec:IntegerProgramming}%
\begin{itemize}[leftmargin=*]%
\item \href{https://www.math.purdue.edu/~yipn/421/2024-10-31-IntLinProg.pdf#page=1}{\textbf{Integer Programming}}: A mathematical optimization problem where the decision variables are restricted to integer values.%
\item \href{https://www.math.purdue.edu/~yipn/421/2024-10-31-IntLinProg.pdf#page=2}{\textbf{Relaxed Problem}}: The linear programming problem obtained by removing the integer constraints from an integer programming problem.%
\item \href{https://www.math.purdue.edu/~yipn/421/2024-10-31-IntLinProg.pdf#page=5}{\textbf{Branch-and-Bound}}: A method for solving integer programming problems by recursively partitioning the feasible region and bounding the objective function value.%
\item \href{https://www.math.purdue.edu/~yipn/421/2024-10-31-IntLinProg.pdf#page=11}{\textbf{Depth-First Search}}: A search strategy for exploring the enumeration tree in Branch and Bound, prioritizing deeper branches before wider ones.%
\item \href{https://www.math.purdue.edu/~yipn/421/2024-10-31-IntLinProg.pdf#page=11}{\textbf{Pruning}}: The process of eliminating branches in the enumeration tree of a Branch and Bound algorithm that cannot lead to a better solution than one already found.%
\item \href{https://www.math.purdue.edu/~yipn/421/2024-10-31-IntLinProg.pdf#page=15}{\textbf{Gomory Cuts}}: Additional constraints added to a linear programming relaxation of an integer program to cut off non-integer solutions.%
\end{itemize}

%
\subsection{Support Vector Machines (SVM)}%
\label{subsec:SupportVectorMachines(SVM)}%
\begin{itemize}[leftmargin=*]%
\item \href{https://www.math.purdue.edu/~yipn/421/2024-10-22-Regression.pdf#page=5}{\textbf{$L^1$ Regression}}: Minimizes the sum of absolute errors between observed and predicted values.%
\item \href{https://www.math.purdue.edu/~yipn/421/2024-10-22-Regression.pdf#page=6}{\textbf{$L^\textbackslash{}infty$ Regression}}: Minimizes the maximum absolute error between observed and predicted values.%
\item \href{https://www.math.purdue.edu/~yipn/421/2024-10-22-Regression.pdf#page=8}{\textbf{Linear Separability}}: A dataset is linearly separable if there exists a hyperplane that perfectly separates the data points into two classes.%
\item \href{https://www.math.purdue.edu/~yipn/421/2024-10-22-Regression.pdf#page=9}{\textbf{Maximum Margin Classifier}}: A classifier that maximizes the minimum distance from the hyperplane to the closest data points.%
\item \href{https://www.math.purdue.edu/~yipn/421/2024-10-22-Regression.pdf#page=11}{\textbf{Lagrangian Formulation of SVM}}: A method to reformulate the SVM optimization problem using Lagrange multipliers.%
\item \href{https://www.math.purdue.edu/~yipn/421/2024-10-22-Regression.pdf#page=12}{\textbf{Geometric Interpretation of SVM}}: Understanding SVM optimization through the geometric concept of maximizing the margin between classes.%
\end{itemize}

%
\section{Problem Types}%
\label{sec:ProblemTypes}%
\subsection{Linear Programming Fundamentals}%
\label{subsec:LinearProgrammingFundamentals}%
\begin{itemize}[leftmargin=*]%
\item \href{https://www.math.purdue.edu/~yipn/421/2024-09-24-SimplexMatrix.pdf#page=1}{\textbf{Representing Linear Programs in Matrix Notation}}: This problem focuses on expressing linear programming problems using matrices and vectors, facilitating the application of the simplex method.%
\item \href{https://www.math.purdue.edu/~yipn/421/2024-09-24-SimplexMatrix.pdf#page=21}{\textbf{Solving a Linear Program Graphically}}: This problem focuses on solving a linear program using a graphical method, identifying the feasible region and the optimal solution point.%
\item \href{https://www.math.purdue.edu/~yipn/421/2024-10-01-DualGeneralLP.pdf#page=2}{\textbf{Handling Unrestricted Variables in Linear Programming}}: This problem addresses how to handle linear programs where variables are not restricted to be non-negative. A common approach is to reformulate the problem by splitting each unrestricted variable into the difference of two non-negative variables.%
\item \href{https://www.math.purdue.edu/~yipn/421/2024-12-03-Duality.pdf#page=3}{\textbf{Solving Linear Programs Graphically}}: A geometric method for solving linear programs with two variables is presented, involving identifying the feasible region and finding the optimal solution at a vertex of the feasible region.%
\item \href{https://www.math.purdue.edu/~yipn/421/2024-12-03-Duality.pdf#page=22}{\textbf{Representing Linear Programs in Matrix Form}}: A linear program is represented in matrix notation, separating variables into basic and non-basic sets, and expressing the objective function and constraints using these sets.%
\item \href{https://www.math.purdue.edu/~yipn/421/2024-12-03-Duality.pdf#page=29}{\textbf{Formulating and Solving Primal and Dual Problems}}: A primal linear programming problem is presented, and its corresponding dual problem is formulated and solved.%
\end{itemize}

%
\subsection{Simplex Method}%
\label{subsec:SimplexMethod}%
\begin{itemize}[leftmargin=*]%
\item \href{https://www.math.purdue.edu/~yipn/421/2024-08-27-ExSimplex.pdf#page=32}{\textbf{Finding an Initial Feasible Solution}}: The Simplex method requires an initial feasible solution. If the origin is not feasible, an auxiliary problem is introduced to find a feasible solution to the original problem.%
\item \href{https://www.math.purdue.edu/~yipn/421/2024-08-27-ExSimplex.pdf#page=30}{\textbf{Handling Infeasible Origins}}: The origin $($0, 0$)$ may not satisfy all constraints in a linear programming problem. Techniques like introducing an auxiliary problem are used to find a feasible starting point.%
\item \href{https://www.math.purdue.edu/~yipn/421/2024-08-27-ExSimplex.pdf#page=26}{\textbf{Using Dictionaries to Represent Solutions}}: The Simplex method uses dictionaries to represent the current solution and to perform pivot operations efficiently. Each dictionary corresponds to a vertex of the feasible region.%
\item \href{https://www.math.purdue.edu/~yipn/421/2024-09-05-SimplexEfficiency.pdf#page=2}{\textbf{Worst-Case Analysis of the Simplex Method}}: The Klee-Minty example demonstrates that the Simplex method can require an exponential number of iterations in the worst case, depending on the pivot rule used.%
\item \href{https://www.math.purdue.edu/~yipn/421/2024-09-05-SimplexEfficiency.pdf#page=3}{\textbf{Impact of Pivot Rules on Simplex Method Efficiency}}: Different pivot rules (e.g., largest-coefficient, largest-increase) significantly affect the number of iterations required by the Simplex method, with some rules exhibiting exponential worst-case behavior while others perform better in practice.%
\item \href{https://www.math.purdue.edu/~yipn/421/2024-09-05-SimplexEfficiency.pdf#page=9}{\textbf{Empirical Analysis of Simplex Method Performance}}: Empirical studies show that the number of iterations in the Simplex method tends to grow polynomially with the problem size in practice, even though worst-case examples exist.%
\item \href{https://www.math.purdue.edu/~yipn/421/2024-09-05-SimplexEfficiency.pdf#page=10}{\textbf{Theoretical Bounds on Simplex Method Iterations}}: Theoretical analysis of randomized pivot rules provides bounds on the expected number of iterations, which are better than exponential but not polynomial.%
\item \href{https://www.math.purdue.edu/~yipn/421/2024-09-05-SimplexEfficiency.pdf#page=11}{\textbf{Smoothed Complexity of the Simplex Method}}: The smoothed complexity analysis suggests that small random perturbations to the input data can significantly improve the performance of the Simplex method, leading to a polynomial number of iterations with high probability.%
\item \href{https://www.math.purdue.edu/~yipn/421/2024-09-24-SimplexMatrix.pdf#page=4}{\textbf{Partitioning Variables into Basic and Non-Basic Sets}}: This problem involves separating the variables of a linear program into basic and non-basic sets, a crucial step in the simplex algorithm.%
\item \href{https://www.math.purdue.edu/~yipn/421/2024-09-24-SimplexMatrix.pdf#page=8}{\textbf{Expressing the Objective Function and Constraints using Basic and Non-Basic Variables}}: This problem deals with rewriting the objective function and constraints of a linear program in terms of basic and non-basic variables, enabling efficient iterative updates in the simplex method.%
\item \href{https://www.math.purdue.edu/~yipn/421/2024-09-24-SimplexMatrix.pdf#page=14}{\textbf{Utilizing Complementary Slackness in Simplex Iterations}}: This problem explores the relationship between primal and dual variables (complementary slackness) within the simplex method's iterative process.%
\item \href{https://www.math.purdue.edu/~yipn/421/2024-09-24-SimplexMatrix.pdf#page=15}{\textbf{Solving a Linear Program using the Simplex Method in Matrix Form}}: This problem involves applying the simplex method in matrix notation to a specific linear programming problem, demonstrating the iterative process of updating the tableau until optimality is reached.%
\item \href{https://www.math.purdue.edu/~yipn/421/2024-10-01-SensitivityDuality.pdf#page=1}{\textbf{Sensitivity Analysis}}: Determining how changes in the constraint values affect the optimal solution and objective function value.%
\item \href{https://www.math.purdue.edu/~yipn/421/2024-10-01-SensitivityDuality.pdf#page=4}{\textbf{Solving Linear Programs using the Simplex Method in Matrix Form}}: Utilizing matrix operations (e.g., inverse of the basis matrix) to solve linear programs using the simplex method.%
\item \href{https://www.math.purdue.edu/~yipn/421/2024-12-03-Duality.pdf#page=11}{\textbf{Performing Sensitivity Analysis}}: The impact of changes in the objective function coefficients ($c$) and constraint values ($b$) on the optimal solution is analyzed using formulas involving the optimal primal and dual solutions.%
\item \href{https://www.math.purdue.edu/~yipn/421/2024-12-03-Duality.pdf#page=16}{\textbf{Determining the Range of Parameter Changes Maintaining Optimality}}: Given an optimal solution to a linear program, the range of values for a parameter in either the objective function or constraints is determined such that the current optimal solution remains optimal.%
\item \href{https://www.math.purdue.edu/~yipn/421/2024-12-03-Duality.pdf#page=26}{\textbf{Applying Sensitivity Analysis to a Specific Problem}}: Sensitivity analysis is applied to a specific linear programming problem to determine the impact of changes in the constraint values on the optimal solution value.%
\end{itemize}

%
\subsection{Duality Theory}%
\label{subsec:DualityTheory}%
\begin{itemize}[leftmargin=*]%
\item \href{https://www.math.purdue.edu/~yipn/421/2024-09-12-Duality.pdf#page=21}{\textbf{Verifying Optimality using Complementary Slackness}}: This problem focuses on using the complementary slackness conditions ($x_j z_j = 0$ and $w_i y_i = 0$) to verify whether a given primal and dual feasible solution pair is optimal.%
\item \href{https://www.math.purdue.edu/~yipn/421/2024-09-12-Duality.pdf#page=6}{\textbf{Demonstrating the Negative Transpose Property in Simplex Iterations}}: This problem involves showing that the negative transpose relationship between the primal and dual tableaux is maintained throughout the simplex method iterations.%
\item \href{https://www.math.purdue.edu/~yipn/421/2024-09-12-Duality.pdf#page=11}{\textbf{Proving the Strong Duality Theorem}}: This problem involves proving that if the primal problem has an optimal solution, then the dual problem also has an optimal solution, and the optimal objective function values are equal. The proof utilizes the structure of the simplex tableau at optimality and the complementary slackness conditions.%
\item \href{https://www.math.purdue.edu/~yipn/421/2024-09-12-Duality.pdf#page=22}{\textbf{Relating Primal and Dual Objective Functions}}: This problem involves demonstrating the relationship between the primal and dual objective functions, showing how the primal objective function is always less than or equal to the dual objective function, and proving their equality at optimality.%
\item \href{https://www.math.purdue.edu/~yipn/421/2024-09-17-DualityGeneral.pdf#page=2}{\textbf{Formulating the Lagrangian for Equality Constraints}}: This problem involves constructing the Lagrangian function $L(x, \textbackslash{}mu) = f_0(x) - \textbackslash{}sum_\textbackslash{}{j=1\textbackslash{}}^n \textbackslash{}mu_j g_j(x)$ to incorporate equality constraints $g_j(x) = 0$ into the optimization problem.%
\item \href{https://www.math.purdue.edu/~yipn/421/2024-09-17-DualityGeneral.pdf#page=7}{\textbf{Formulating the Lagrangian for Inequality and Equality Constraints}}: This problem involves constructing the Lagrangian function $L(x, \textbackslash{}lambda, \textbackslash{}mu) = f_0(x) - \textbackslash{}sum_\textbackslash{}{i=1\textbackslash{}}^m \textbackslash{}lambda_i f_i(x) - \textbackslash{}sum_\textbackslash{}{j=1\textbackslash{}}^n \textbackslash{}mu_j g_j(x)$ to incorporate both inequality constraints $f_i(x) \textbackslash{}le 0$ and equality constraints $g_j(x) = 0$ into the optimization problem.%
\item \href{https://www.math.purdue.edu/~yipn/421/2024-09-17-DualityGeneral.pdf#page=5}{\textbf{Establishing Weak Duality}}: This problem involves proving that the optimal value of the primal problem is less than or equal to the optimal value of the dual problem, i.e., $\textbackslash{}max_\textbackslash{}{x \textbackslash{}in \textbackslash{}mathcal\textbackslash{}{G\textbackslash{}}\textbackslash{}} \textbackslash{}xi(x) \textbackslash{}le \textbackslash{}min_\textbackslash{}{\textbackslash{}lambda \textbackslash{}ge 0, \textbackslash{}mu\textbackslash{}} \textbackslash{}xi(\textbackslash{}lambda, \textbackslash{}mu)$.%
\item \href{https://www.math.purdue.edu/~yipn/421/2024-09-17-DualityGeneral.pdf#page=17}{\textbf{Establishing Strong Duality}}: This problem involves proving that under certain conditions (e.g., the primal problem having a finite optimal solution), the optimal values of the primal and dual problems are equal, i.e., $\textbackslash{}xi(x^*) = \textbackslash{}xi(\textbackslash{}lambda^)$.%
\item \href{https://www.math.purdue.edu/~yipn/421/2024-09-17-DualityGeneral.pdf#page=15}{\textbf{Deriving the Dual Problem from the Lagrangian}}: This problem involves deriving the dual problem by analyzing the maximum of the Lagrangian function with respect to the primal variables, leading to the dual objective function and constraints.%
\item \href{https://www.math.purdue.edu/~yipn/421/2024-09-17-DualityGeneral.pdf#page=22}{\textbf{Applying Complementary Slackness}}: This problem involves using the complementary slackness condition ($x^*_i \textbackslash{}lambda^_i = 0$ for all $i$) to verify the optimality of solutions to the primal and dual problems.%
\item \href{https://www.math.purdue.edu/~yipn/421/2024-10-01-DualGeneralLP.pdf#page=1}{\textbf{Deriving the Dual Problem from the Primal Problem}}: This problem focuses on deriving the dual linear program from a given primal linear program. The derivation involves techniques such as forming the Lagrangian and manipulating inequalities to obtain the dual objective function and constraints.%
\item \href{https://www.math.purdue.edu/~yipn/421/2024-10-01-SensitivityDuality.pdf#page=2}{\textbf{Deriving the Dual Problem}}: Formulating the dual linear program from a given primal linear program.%
\item \href{https://www.math.purdue.edu/~yipn/421/2024-12-03-Duality.pdf#page=5}{\textbf{Formulating the Dual Problem}}: Given a primal linear program, the corresponding dual linear program is formulated, showing the transformation of coefficients and constraints.%
\item \href{https://www.math.purdue.edu/~yipn/421/2024-12-03-Duality.pdf#page=8}{\textbf{Finding Optimal Solutions Using Duality}}: The optimal solutions to both the primal and dual problems are found and verified using complementary slackness conditions.%
\item \href{https://www.math.purdue.edu/~yipn/421/2024-12-03-Duality.pdf#page=9}{\textbf{Geometric Interpretation of Dual Constraints}}: The dual constraints are interpreted geometrically as half-spaces in the primal variable space, and the optimal dual solution is identified as the intersection of these half-spaces.%
\item \href{https://www.math.purdue.edu/~yipn/421/2024-12-03-Duality.pdf#page=31}{\textbf{Analyzing Complementary Slackness Conditions}}: The complementary slackness conditions are analyzed to determine the optimal solution and the range of parameter values that maintain optimality.%
\end{itemize}

%
\subsection{Advanced Linear Programming Topics}%
\label{subsec:AdvancedLinearProgrammingTopics}%
\begin{itemize}[leftmargin=*]%
\item \href{https://www.math.purdue.edu/~yipn/421/2024-09-05-SimplexEfficiency.pdf#page=12}{\textbf{Existence of Polynomial-Time Algorithms for Linear Programming}}: Polynomial-time algorithms for linear programming exist, such as interior-point methods, which contrast with the Simplex method's potential for exponential worst-case behavior.%
\item \href{https://www.math.purdue.edu/~yipn/421/2024-10-28-ApplsLinProg.pdf#page=30}{\textbf{Sparse Solutions of Linear Systems}}: Finding a solution to an underdetermined system of linear equations that has a limited number of non-zero elements.%
\end{itemize}

%
\subsection{Network Flow Problems}%
\label{subsec:NetworkFlowProblems}%
\begin{itemize}[leftmargin=*]%
\item \href{https://www.math.purdue.edu/~yipn/421/2024-10-28-ApplsLinProg.pdf#page=19}{\textbf{Maximum Weight Matching}}: Finding a matching in a weighted graph that maximizes the total weight of the selected edges.%
\item \href{https://www.math.purdue.edu/~yipn/421/2024-11-05-NetFlows.pdf#page=2}{\textbf{Formulating Network Flow Problems as Linear Programs}}: This problem involves representing a network flow problem, characterized by nodes, arcs, supply/demand at nodes, and arc costs, as a linear program with an objective function minimizing total transportation cost and constraints ensuring flow conservation at each node.%
\item \href{https://www.math.purdue.edu/~yipn/421/2024-11-05-NetFlows.pdf#page=6}{\textbf{Analyzing the Incidence Matrix of a Network}}: This problem focuses on understanding the properties of the incidence matrix, including its rank and the relationship between its submatrices and spanning trees in the network. The goal is to understand how the structure of the incidence matrix relates to the solvability and solution methods for the network flow problem.%
\item \href{https://www.math.purdue.edu/~yipn/421/2024-11-05-NetFlows.pdf#page=10}{\textbf{Relating Spanning Trees to Bases in Network Flow Problems}}: This problem explores the connection between spanning trees in a network and bases in the simplex method. The goal is to understand how spanning trees can be used to represent basic feasible solutions and how this relationship can be exploited to solve network flow problems efficiently.%
\item \href{https://www.math.purdue.edu/~yipn/421/2024-11-05-NetFlows.pdf#page=22}{\textbf{Applying the Primal Simplex Method to Network Flow Problems}}: This problem involves using the primal simplex method to iteratively improve a feasible solution to a network flow problem, represented by a spanning tree, until an optimal solution is reached. The method involves identifying entering and leaving variables, updating the spanning tree, and updating primal and dual variables.%
\item \href{https://www.math.purdue.edu/~yipn/421/2024-11-07-NetSimplex.pdf#page=2}{\textbf{Finding a primal feasible flow in a network}}: This problem focuses on finding a flow vector $x$ such that $\textbackslash{}tilde\textbackslash{}{B\textbackslash{}}\textbackslash{}tilde\textbackslash{}{x\textbackslash{}} = \textbackslash{}tilde\textbackslash{}{b\textbackslash{}}$ and $x_\textbackslash{}{ij\textbackslash{}} = 0$ for $(i,j) \textbackslash{}notin T$, where $T$ is a spanning tree. The solution might have negative components, requiring further adjustments.%
\item \href{https://www.math.purdue.edu/~yipn/421/2024-11-07-NetSimplex.pdf#page=6}{\textbf{Selecting entering and leaving arcs in the primal network simplex method}}: This problem focuses on choosing an entering arc with $z_\textbackslash{}{ij\textbackslash{}} < 0$ and a leaving arc to form a cycle, updating primal and dual flows to improve the solution.%
\item \href{https://www.math.purdue.edu/~yipn/421/2024-11-07-NetSimplex.pdf#page=35}{\textbf{Selecting the entering arc in the dual network simplex method}}: This problem involves choosing an entering arc that bridges two subtrees in the opposite direction from the leaving arc, and has the smallest dual slack among all such arcs.%
\item \href{https://www.math.purdue.edu/~yipn/421/2024-11-07-NetSimplex.pdf#page=6}{\textbf{Updating primal and dual flows after selecting entering and leaving arcs}}: This problem involves updating the primal and dual flows along the cycle formed by the entering and leaving arcs, ensuring that the spanning tree remains valid and the flow values are adjusted correctly.%
\item \href{https://www.math.purdue.edu/~yipn/421/2024-11-07-NetSimplex.pdf#page=14}{\textbf{Updating dual variables in the primal network simplex method}}: This problem focuses on updating dual variables after a primal pivot step, considering the cases where the entering arc crosses between trees containing and not containing the root node.%
\item \href{https://www.math.purdue.edu/~yipn/421/2024-11-07-NetSimplex.pdf#page=45}{\textbf{Finding an initial feasible solution using a two-phased network simplex method}}: This problem involves modifying costs to ensure dual feasibility of an initial tree solution, then using the primal network simplex method to find an optimal solution.%
\item \href{https://www.math.purdue.edu/~yipn/421/2024-11-07-NetSimplex.pdf#page=45}{\textbf{Finding an initial feasible solution using the primal network simplex method}}: This problem involves modifying supplies to ensure primal feasibility of an initial tree solution, then using the dual network simplex method to find an optimal solution.%
\item \href{https://www.math.purdue.edu/~yipn/421/2024-11-07-NetSimplex.pdf#page=45}{\textbf{Solving network flow problems using a parametric self-dual method}}: This problem involves interweaving primal and dual pivots to reduce a perturbation parameter $\textbackslash{}mu$ from $\textbackslash{}infty$ to zero.%
\item \href{https://www.math.purdue.edu/~yipn/421/2024-11-19-MaxFlowMinCut.pdf#page=1}{\textbf{Maximum Flow Problem}}: Finding the maximum flow from a source node to a sink node in a network, subject to capacity constraints on the edges.%
\item \href{https://www.math.purdue.edu/~yipn/421/2024-11-19-MaxFlowMinCut.pdf#page=8}{\textbf{Max-Flow Min-Cut Theorem}}: Proving that the maximum flow in a network is equal to the minimum capacity of any cut separating the source and sink nodes.%
\item \href{https://www.math.purdue.edu/~yipn/421/2024-11-19-MaxFlowMinCut.pdf#page=10}{\textbf{Implementing the Ford-Fulkerson Algorithm}}: Describing and implementing the steps of the Ford-Fulkerson algorithm, which iteratively finds augmenting paths to increase the flow until the maximum flow is reached.%
\item \href{https://www.math.purdue.edu/~yipn/421/2024-11-19-MaxFlowMinCut.pdf#page=14}{\textbf{Linear Programming Formulation of Maximum Flow}}: Formulating the maximum flow problem as a linear program, including the objective function and constraints.%
\end{itemize}

%
\subsection{Integer Programming}%
\label{subsec:IntegerProgramming}%
\begin{itemize}[leftmargin=*]%
\item \href{https://www.math.purdue.edu/~yipn/421/2024-10-28-ApplsLinProg.pdf#page=24}{\textbf{Knapsack Problem}}: Selecting a subset of items with weights and values to maximize the total value while not exceeding a weight capacity.%
\item \href{https://www.math.purdue.edu/~yipn/421/2024-10-31-IntLinProg.pdf#page=1}{\textbf{Integer Programming}}: A linear programming problem where some or all variables are restricted to integer values.%
\item \href{https://www.math.purdue.edu/~yipn/421/2024-10-31-IntLinProg.pdf#page=5}{\textbf{Solving Integer Programs using Branch and Bound}}: A method for solving integer programs by recursively partitioning the feasible region and solving relaxed linear programs.%
\item \href{https://www.math.purdue.edu/~yipn/421/2024-10-31-IntLinProg.pdf#page=15}{\textbf{Applying Gomory Cuts}}: A method to generate additional constraints (cuts) to eliminate non-integer solutions from the feasible region of an integer program.%
\end{itemize}

%
\subsection{Applications of Linear Programming}%
\label{subsec:ApplicationsofLinearProgramming}%
\begin{itemize}[leftmargin=*]%
\item \href{https://www.math.purdue.edu/~yipn/421/2024-09-19-DualityExamples.pdf#page=5}{\textbf{Formulating the Diet Problem as a Linear Program}}: This problem involves translating the description of a diet problem (minimizing cost while meeting nutritional requirements) into a mathematical linear program with an objective function and constraints.%
\item \href{https://www.math.purdue.edu/~yipn/421/2024-09-19-DualityExamples.pdf#page=8}{\textbf{Formulating the Dual of the Diet Problem}}: This problem involves deriving the dual linear program from the primal linear program representing the diet problem, interpreting the dual variables in the context of the problem.%
\item \href{https://www.math.purdue.edu/~yipn/421/2024-09-19-DualityExamples.pdf#page=3}{\textbf{Representing the Diet Problem using Tables}}: This problem involves representing the data of the diet problem (nutrient content of foods, minimum nutrient requirements, and food costs) in a tabular format to facilitate the formulation of the linear program.%
\item \href{https://www.math.purdue.edu/~yipn/421/2024-09-19-DualityExamples.pdf#page=7}{\textbf{Transforming the Diet Problem into a Pill-Selling Problem}}: This problem involves reinterpreting the diet problem as a profit maximization problem by considering the sale of nutrient pills instead of food consumption.%
\item \href{https://www.math.purdue.edu/~yipn/421/2024-10-28-ApplsLinProg.pdf#page=3}{\textbf{Production Planning}}: Formulating and solving a linear program to optimize ice cream production over a year, minimizing costs associated with changing production levels and storage, given predicted monthly demands.%
\item \href{https://www.math.purdue.edu/~yipn/421/2024-10-28-ApplsLinProg.pdf#page=7}{\textbf{Workforce Planning}}: Determining the optimal mix of regular and temporary employees to minimize labor costs while meeting fluctuating monthly labor demands.%
\item \href{https://www.math.purdue.edu/~yipn/421/2024-10-28-ApplsLinProg.pdf#page=22}{\textbf{Machine Scheduling}}: Optimizing the scheduling of jobs on multiple machines to minimize the total completion time.%
\item \href{https://www.math.purdue.edu/~yipn/421/2024-10-28-ApplsLinProg.pdf#page=25}{\textbf{Cutting Stock Problem}}: Determining the optimal way to cut large rolls of material into smaller rolls of various widths to minimize waste.%
\end{itemize}

%
\subsection{Geometric and Algorithmic Aspects}%
\label{subsec:GeometricandAlgorithmicAspects}%
\begin{itemize}[leftmargin=*]%
\item \href{https://www.math.purdue.edu/~yipn/421/2024-08-27-ExSimplex.pdf#page=17}{\textbf{Optimizing in Higher Dimensions}}: The Simplex method is extended to linear programming problems with more than two variables, requiring iterative pivoting operations to find the optimal solution.%
\item \href{https://www.math.purdue.edu/~yipn/421/2024-10-22-Regression.pdf#page=15}{\textbf{Maximum Inscribed Circle in a Convex Polyhedron}}: This problem involves finding the largest circle that can be inscribed within a given convex polyhedron.%
\item \href{https://www.math.purdue.edu/~yipn/421/2024-10-22-Regression.pdf#page=22}{\textbf{Smallest Ball Problem}}: This problem aims to find the smallest radius ball that encloses a given set of points in a d-dimensional space.%
\item \href{https://www.math.purdue.edu/~yipn/421/2024-10-28-ApplsLinProg.pdf#page=17}{\textbf{Smallest Enclosing Ball Problem}}: Finding the smallest ball (hypersphere) that contains a given set of points in a $d$-dimensional space.%
\end{itemize}

%
\subsection{Machine Learning Connections}%
\label{subsec:MachineLearningConnections}%
\begin{itemize}[leftmargin=*]%
\item \href{https://www.math.purdue.edu/~yipn/421/2024-10-22-Regression.pdf#page=1}{\textbf{Regression}}: This problem involves finding parameters that best fit a linear model to a set of data points, minimizing different error functions such as $L^2$, $L^1$, and $L^\textbackslash{}infty$ norms.%
\item \href{https://www.math.purdue.edu/~yipn/421/2024-10-22-Regression.pdf#page=7}{\textbf{Binary Classification}}: This problem focuses on finding a hyperplane that separates data points into two classes.%
\item \href{https://www.math.purdue.edu/~yipn/421/2024-10-22-Regression.pdf#page=9}{\textbf{Maximum Margin Classifier}}: This problem aims to find the hyperplane that maximizes the minimum distance to the closest data points of either class.%
\end{itemize}

%
\section{Algorithm Solutions}%
\label{sec:AlgorithmSolutions}%
\subsection{Linear Programming Fundamentals}%
\label{subsec:LinearProgrammingFundamentals}%
\begin{itemize}[leftmargin=*]%
\item \href{https://www.math.purdue.edu/~yipn/421/2024-09-12-Duality.pdf#page=2}{\textbf{Weak Duality}}: $c^T x \textbackslash{}le b^T y$%
\item \href{https://www.math.purdue.edu/~yipn/421/2024-09-12-Duality.pdf#page=11}{\textbf{Strong Duality}}: If (P) has an optimal solution $x^*$, then (D) has an optimal solution $y^*$, and $c^T x^* = b^T y^*$.%
\item \href{https://www.math.purdue.edu/~yipn/421/2024-09-12-Duality.pdf#page=21}{\textbf{Complementary Slackness}}: Primal feasible $x$ and dual feasible $y$ are optimal if and only if $x_j z_j = 0$ for all $j$ and $w_i y_i = 0$ for all $i$.%
\item \href{https://www.math.purdue.edu/~yipn/421/2024-09-17-DualityGeneral.pdf#page=7}{\textbf{Lagrangian Function}}: $L(x,\textbackslash{}lambda,\textbackslash{}mu) = f_0(x) - \textbackslash{}sum_\textbackslash{}{i=1\textbackslash{}}^m \textbackslash{}lambda_i f_i(x) - \textbackslash{}sum_\textbackslash{}{j=1\textbackslash{}}^n \textbackslash{}mu_j g_j(x)$%
\item \href{https://www.math.purdue.edu/~yipn/421/2024-10-01-DualGeneralLP.pdf#page=4}{\textbf{Dual Problem Derivation using Lagrangian Multipliers}}: The dual problem is derived from the primal problem by constructing the Lagrangian function, $\textbackslash{}mathcal\textbackslash{}{L\textbackslash{}}(\textbackslash{}vec\textbackslash{}{x\textbackslash{}}, \textbackslash{}vec\textbackslash{}{y\textbackslash{}}, \textbackslash{}vec\textbackslash{}{z\textbackslash{}}) = \textbackslash{}vec\textbackslash{}{c\textbackslash{}}^T \textbackslash{}vec\textbackslash{}{x\textbackslash{}} - \textbackslash{}vec\textbackslash{}{y\textbackslash{}}^T (\textbackslash{}vec\textbackslash{}{A\textbackslash{}} \textbackslash{}vec\textbackslash{}{x\textbackslash{}} - \textbackslash{}vec\textbackslash{}{b\textbackslash{}}) - \textbackslash{}vec\textbackslash{}{z\textbackslash{}}^T (\textbackslash{}vec\textbackslash{}{x\textbackslash{}} - \textbackslash{}vec\textbackslash{}{u\textbackslash{}})$, where $\textbackslash{}vec\textbackslash{}{y\textbackslash{}}$ and $\textbackslash{}vec\textbackslash{}{z\textbackslash{}}$ are Lagrange multipliers. By taking the partial derivative with respect to $\textbackslash{}vec\textbackslash{}{x\textbackslash{}}$ and setting it to zero, and considering the constraints $\textbackslash{}vec\textbackslash{}{y\textbackslash{}} \textbackslash{}ge \textbackslash{}vec\textbackslash{}{0\textbackslash{}}$ and $\textbackslash{}vec\textbackslash{}{z\textbackslash{}} \textbackslash{}ge \textbackslash{}vec\textbackslash{}{0\textbackslash{}}$, the dual problem is obtained.%
\item \href{https://www.math.purdue.edu/~yipn/421/2024-10-01-DualGeneralLP.pdf#page=2}{\textbf{Reformulation of Primal Problem with Unrestricted Variables}}: A primal problem with unrestricted variables $\textbackslash{}vec\textbackslash{}{x\textbackslash{}}$ can be reformulated by splitting each variable into its positive and negative parts, $\textbackslash{}vec\textbackslash{}{x\textbackslash{}} = \textbackslash{}vec\textbackslash{}{x\textbackslash{}}^+ - \textbackslash{}vec\textbackslash{}{x\textbackslash{}}^-$, where $\textbackslash{}vec\textbackslash{}{x\textbackslash{}}^+ \textbackslash{}ge \textbackslash{}vec\textbackslash{}{0\textbackslash{}}$ and $\textbackslash{}vec\textbackslash{}{x\textbackslash{}}^- \textbackslash{}ge \textbackslash{}vec\textbackslash{}{0\textbackslash{}}$. This allows the application of standard linear programming techniques.%
\item \href{https://www.math.purdue.edu/~yipn/421/2024-10-01-SensitivityDuality.pdf#page=7}{\textbf{Sensitivity Analysis}}: This technique analyzes how changes in the parameters of a linear programming problem (e.g., objective function coefficients, constraint coefficients, or right-hand side values) affect the optimal solution and objective function value. It often involves calculating the range of parameter changes that maintain the current optimal basis.%
\item \href{https://www.math.purdue.edu/~yipn/421/2024-10-01-SensitivityDuality.pdf#page=10}{\textbf{Shadow Price}}: The shadow price of a constraint is the rate of change in the optimal objective function value for a unit change in the right-hand side of that constraint. It is given by the corresponding optimal dual variable.%
\item \href{https://www.math.purdue.edu/~yipn/421/2024-10-22-Convexity.pdf#page=8}{\textbf{Caratheodory's Theorem}}: If $z \textbackslash{}in \textbackslash{}text\textbackslash{}{Conv\textbackslash{}}(S)$, then $z$ can be written as a convex combination of at most $m+1$ points from $S$.%
\item \href{https://www.math.purdue.edu/~yipn/421/2024-10-22-Convexity.pdf#page=13}{\textbf{Farkas' Lemma}}: The system $Ax \textbackslash{}le b$ has no solution if and only if there exists a vector $y \textbackslash{}ge 0$ such that $y^*T A = 0$ and $y^*T b < 0$.%
\item \href{https://www.math.purdue.edu/~yipn/421/2024-10-22-Convexity.pdf#page=10}{\textbf{Separation Theorem}}: Let $P$ and $\textbackslash{}bar\textbackslash{}{P\textbackslash{}}$ be two disjoint nonempty polyhedra in $\textbackslash{}mathbb\textbackslash{}{R\textbackslash{}}^n$. Then there exist disjoint half-spaces $H$ and $\textbackslash{}bar\textbackslash{}{H\textbackslash{}}$ such that $P \textbackslash{}subset H$ and $\textbackslash{}bar\textbackslash{}{P\textbackslash{}} \textbackslash{}subset \textbackslash{}bar\textbackslash{}{H\textbackslash{}}$.%
\item \href{https://www.math.purdue.edu/~yipn/421/2024-12-03-Duality.pdf#page=3}{\textbf{Graphical Method}}: Solve a linear program with two variables by plotting the constraints to define the feasible region and then finding the vertex that maximizes (or minimizes) the objective function.%
\item \href{https://www.math.purdue.edu/~yipn/421/2024-12-03-Duality.pdf#page=1}{\textbf{Dual Problem Derivation}}: Given a primal problem $\textbackslash{}text\textbackslash{}{max \textbackslash{}} c^T x \textbackslash{}text\textbackslash{}{ s.t. \textbackslash{}} Ax \textbackslash{}le b, x \textbackslash{}ge 0$, the dual problem is $\textbackslash{}text\textbackslash{}{min \textbackslash{}} b^T y \textbackslash{}text\textbackslash{}{ s.t. \textbackslash{}} A^T y \textbackslash{}ge c, y \textbackslash{}ge 0$.%
\end{itemize}

%
\subsection{Simplex Method and Variants}%
\label{subsec:SimplexMethodandVariants}%
\begin{itemize}[leftmargin=*]%
\item \href{https://www.math.purdue.edu/~yipn/421/2024-08-27-ExSimplex.pdf#page=6}{\textbf{Simplex Method}}: Iteratively move from one vertex of the feasible region to another by setting non-basic variables to zero and choosing a new set of (non)basic variables to improve the objective function value. This is equivalent to moving from vertex to vertex in the feasible region.%
\item \href{https://www.math.purdue.edu/~yipn/421/2024-08-27-ExSimplex.pdf#page=32}{\textbf{Auxiliary Problem}}: When the origin is not feasible, introduce a new variable $x_0$ and modify the constraints to create an auxiliary problem that maximizes $-x_0$. Solve this auxiliary problem using the Simplex method. If the optimal solution has $x_0 = 0$, then the solution is feasible for the original problem.%
\item \href{https://www.math.purdue.edu/~yipn/421/2024-09-05-SimplexEfficiency.pdf#page=4}{\textbf{Largest-Coefficient Rule}}: Choose the variable with the largest coefficient in the objective function row to enter the basis.%
\item \href{https://www.math.purdue.edu/~yipn/421/2024-09-05-SimplexEfficiency.pdf#page=6}{\textbf{Largest-Increase Rule}}: Choose the variable whose entrance into the basis brings about the largest increase in the objective function.%
\item \href{https://www.math.purdue.edu/~yipn/421/2024-09-05-SimplexEfficiency.pdf#page=10}{\textbf{Bland's Rule}}: Choose the variable with the smallest index among those eligible to enter the basis.%
\item \href{https://www.math.purdue.edu/~yipn/421/2024-09-05-SimplexEfficiency.pdf#page=10}{\textbf{Randomized Pivot Rule}}: Randomly permute the indices of the variables; then use Bland's rule for entering and the lexicographic method for leaving variables.%
\item \href{https://www.math.purdue.edu/~yipn/421/2024-09-24-SimplexMatrix.pdf#page=8}{\textbf{Simplex Method in Matrix Form}}: The simplex method iteratively improves a feasible solution by exchanging basic and non-basic variables until an optimal solution is found. This involves updating the objective function and constraints using matrix operations, specifically involving the basis matrix $B$ and its inverse $B^\textbackslash{}{-1\textbackslash{}}$, and the non-basic variable matrix $N$. The objective function is expressed as $\textbackslash{}mathcal\textbackslash{}{Z\textbackslash{}} = \textbackslash{}mathbf\textbackslash{}{c\textbackslash{}}_B^T B^\textbackslash{}{-1\textbackslash{}} \textbackslash{}mathbf\textbackslash{}{b\textbackslash{}} + (\textbackslash{}mathbf\textbackslash{}{c\textbackslash{}}_N^T - \textbackslash{}mathbf\textbackslash{}{c\textbackslash{}}_B^T B^\textbackslash{}{-1\textbackslash{}} N) \textbackslash{}mathbf\textbackslash{}{x\textbackslash{}}_N$, where $\textbackslash{}mathbf\textbackslash{}{c\textbackslash{}}_B$ and $\textbackslash{}mathbf\textbackslash{}{c\textbackslash{}}_N$ are the cost vectors for basic and non-basic variables, respectively, and $\textbackslash{}mathbf\textbackslash{}{b\textbackslash{}}$ is the right-hand side vector of the constraints.%
\end{itemize}

%
\subsection{Regression and Optimization}%
\label{subsec:RegressionandOptimization}%
\begin{itemize}[leftmargin=*]%
\item \href{https://www.math.purdue.edu/~yipn/421/2024-10-22-Regression.pdf#page=5}{\textbf{$L^1$ Regression}}: $\textbackslash{}min_\textbackslash{}{x_1, \textbackslash{}dots, x_n\textbackslash{}} \textbackslash{}sum_\textbackslash{}{i=1\textbackslash{}}^m |b_i - \textbackslash{}sum_\textbackslash{}{j=1\textbackslash{}}^n a_\textbackslash{}{ij\textbackslash{}} x_j|$ or $\textbackslash{}min_X \textbackslash{}|b - AX\textbackslash{}|_1$, solved by introducing slack variables $\textbackslash{}hat\textbackslash{}{\textbackslash{}varepsilon\textbackslash{}}_i$ such that $-\textbackslash{}hat\textbackslash{}{\textbackslash{}varepsilon\textbackslash{}}_i \textbackslash{}le b_i - \textbackslash{}sum_\textbackslash{}{j=1\textbackslash{}}^n a_\textbackslash{}{ij\textbackslash{}} x_j \textbackslash{}le \textbackslash{}hat\textbackslash{}{\textbackslash{}varepsilon\textbackslash{}}_i$ and minimizing $\textbackslash{}sum_\textbackslash{}{i=1\textbackslash{}}^m \textbackslash{}hat\textbackslash{}{\textbackslash{}varepsilon\textbackslash{}}_i$.%
\item \href{https://www.math.purdue.edu/~yipn/421/2024-10-22-Regression.pdf#page=6}{\textbf{$L^\textbackslash{}infty$ Regression}}: $\textbackslash{}min_\textbackslash{}{x_1, \textbackslash{}dots, x_n\textbackslash{}} \textbackslash{}max_\textbackslash{}{i\textbackslash{}} |b_i - \textbackslash{}sum_\textbackslash{}{j=1\textbackslash{}}^n a_\textbackslash{}{ij\textbackslash{}} x_j|$ or $\textbackslash{}min_X \textbackslash{}|b - AX\textbackslash{}|_\textbackslash{}infty$, solved by introducing a slack variable $\textbackslash{}delta$ such that $-\textbackslash{}delta \textbackslash{}le b_i - \textbackslash{}sum_\textbackslash{}{j=1\textbackslash{}}^n a_\textbackslash{}{ij\textbackslash{}} x_j \textbackslash{}le \textbackslash{}delta, \textbackslash{}quad \textbackslash{}forall i$ and minimizing $\textbackslash{}delta$.%
\item \href{https://www.math.purdue.edu/~yipn/421/2024-10-22-Regression.pdf#page=11}{\textbf{Lagrangian Formulation of SVM}}: $\textbackslash{}min_\textbackslash{}{\textbackslash{}alpha\textbackslash{}} \textbackslash{}frac\textbackslash{}{1\textbackslash{}}\textbackslash{}{2\textbackslash{}} \textbackslash{}sum_\textbackslash{}{i=1\textbackslash{}}^N \textbackslash{}sum_\textbackslash{}{j=1\textbackslash{}}^N \textbackslash{}alpha_i \textbackslash{}alpha_j y_i y_j x_i^T x_j - \textbackslash{}sum_\textbackslash{}{i=1\textbackslash{}}^N \textbackslash{}alpha_i$ subject to $\textbackslash{}sum_\textbackslash{}{i=1\textbackslash{}}^N \textbackslash{}alpha_i y_i = 0$ and $0 \textbackslash{}le \textbackslash{}alpha_i \textbackslash{}le C, \textbackslash{}quad \textbackslash{}forall i$.%
\item \href{https://www.math.purdue.edu/~yipn/421/2024-10-22-Regression.pdf#page=18}{\textbf{Distance of a Point from a Hyperplane}}: $d_i = \textbackslash{}frac\textbackslash{}{|a_1 s_i + a_2 s_i^2 - b'|\textbackslash{}}\textbackslash{}{\textbackslash{}sqrt\textbackslash{}{a_1^2 + a_2^2\textbackslash{}}\textbackslash{}}$%
\item \href{https://www.math.purdue.edu/~yipn/421/2024-10-28-ApplsLinProg.pdf#page=18}{\textbf{Smallest Enclosing Ball}}: minimize $x^*T Q^T Q x - \textbackslash{}sum_\textbackslash{}{j=1\textbackslash{}}^n x_j p_j^T p_j$ subject to $\textbackslash{}sum_\textbackslash{}{j=1\textbackslash{}}^n x_j = 1$, $x_j \textbackslash{}ge 0$%
\item \href{https://www.math.purdue.edu/~yipn/421/2024-10-28-ApplsLinProg.pdf#page=34}{\textbf{Sparse Solutions of Linear System}}: $\textbackslash{}min \textbackslash{}|x\textbackslash{}|_1$ s.t. $Ax = b$%
\end{itemize}

%
\subsection{Integer and Combinatorial Optimization}%
\label{subsec:IntegerandCombinatorialOptimization}%
\begin{itemize}[leftmargin=*]%
\item \href{https://www.math.purdue.edu/~yipn/421/2024-10-28-ApplsLinProg.pdf#page=21}{\textbf{Maximum Weight Matching}}: $\textbackslash{}max \textbackslash{}sum_\textbackslash{}{e \textbackslash{}in E\textbackslash{}} w_e x_e$ s.t. $\textbackslash{}sum_\textbackslash{}{e \textbackslash{}in \textbackslash{}delta(v)\textbackslash{}} x_e = 1$ for all vertices $v$, $x_e \textbackslash{}in \textbackslash{}\textbackslash{}{$0, 1$\textbackslash{}\textbackslash{}}$ for all edges $e$%
\item \href{https://www.math.purdue.edu/~yipn/421/2024-10-28-ApplsLinProg.pdf#page=23}{\textbf{Machine Scheduling}}: $\textbackslash{}min \textbackslash{}pi$ s.t. $\textbackslash{}sum_\textbackslash{}{i \textbackslash{}in M\textbackslash{}} x_\textbackslash{}{ij\textbackslash{}} = 1$ for $j \textbackslash{}in J$, $\textbackslash{}sum_\textbackslash{}{j \textbackslash{}in J\textbackslash{}} d_\textbackslash{}{ij\textbackslash{}} x_\textbackslash{}{ij\textbackslash{}} \textbackslash{}le \textbackslash{}pi$ for $i \textbackslash{}in M$, $x_\textbackslash{}{ij\textbackslash{}} \textbackslash{}in \textbackslash{}\textbackslash{}{$0, 1$\textbackslash{}\textbackslash{}}$ for all $i, j$%
\item \href{https://www.math.purdue.edu/~yipn/421/2024-10-28-ApplsLinProg.pdf#page=24}{\textbf{0-1 Knapsack Problem}}: maximize $\textbackslash{}sum_\textbackslash{}{j=1\textbackslash{}}^n c_j x_j$ subject to $\textbackslash{}sum_\textbackslash{}{j=1\textbackslash{}}^n a_j x_j \textbackslash{}le b$, $x_j \textbackslash{}in \textbackslash{}\textbackslash{}{$0, 1$\textbackslash{}\textbackslash{}}$, $j = $1, 2$, \textbackslash{}dots, n$%
\item \href{https://www.math.purdue.edu/~yipn/421/2024-10-28-ApplsLinProg.pdf#page=29}{\textbf{Cutting Stock Problem}}: $\textbackslash{}min \textbackslash{}sum_\textbackslash{}{i=1\textbackslash{}}^\textbackslash{}{12\textbackslash{}} x_i$ subject to constraints based on cutting patterns $P_i$ and demand for each roll width%
\item \href{https://www.math.purdue.edu/~yipn/421/2024-10-31-IntLinProg.pdf#page=5}{\textbf{Branch and Bound}}: This algorithm solves integer programming problems by recursively partitioning the feasible region into subproblems, solving the relaxed linear program for each subproblem, and using the objective function values to prune branches that cannot lead to a better integer solution.%
\item \href{https://www.math.purdue.edu/~yipn/421/2024-10-31-IntLinProg.pdf#page=15}{\textbf{Gomory Cuts}}: This algorithm adds cutting planes (constraints) to the relaxed linear program to eliminate non-integer solutions. The cuts are derived from the fractional parts of the coefficients in the optimal dictionary of the relaxed problem. The algorithm iteratively adds cuts until an integer solution is found.%
\end{itemize}

%
\subsection{Network Flows}%
\label{subsec:NetworkFlows}%
\begin{itemize}[leftmargin=*]%
\item \href{https://www.math.purdue.edu/~yipn/421/2024-11-05-NetFlows.pdf#page=2}{\textbf{Network Flow Problem Formulation}}: The network flow problem is formulated as a linear program with the objective function $\textbackslash{}text\textbackslash{}{min\textbackslash{}} \textbackslash{}sum_\textbackslash{}{(i,j) \textbackslash{}in \textbackslash{}mathcal\textbackslash{}{A\textbackslash{}}\textbackslash{}} C_\textbackslash{}{ij\textbackslash{}} X_\textbackslash{}{ij\textbackslash{}}$ subject to flow conservation constraints $\textbackslash{}sum_\textbackslash{}{i\textbackslash{}} x_\textbackslash{}{ik\textbackslash{}} - \textbackslash{}sum_\textbackslash{}{j\textbackslash{}} x_\textbackslash{}{kj\textbackslash{}} = -b_k$ at each node $k$ and non-negativity constraints $X \textbackslash{}ge 0$.%
\item \href{https://www.math.purdue.edu/~yipn/421/2024-11-05-NetFlows.pdf#page=22}{\textbf{Primal Simplex Method for Network Flows}}: The primal simplex method is adapted for network flows by iteratively selecting an entering variable (arc with the most negative reduced cost), identifying a leaving variable (arc in the cycle formed by the entering variable and the spanning tree), updating the spanning tree, and updating primal and dual variables until all dual slacks are non-negative.%
\item \href{https://www.math.purdue.edu/~yipn/421/2024-11-05-NetFlows.pdf#page=15}{\textbf{Finding a Spanning Tree Basis}}: A square submatrix of the modified incidence matrix $\textbackslash{}tilde\textbackslash{}{A\textbackslash{}}$ (incidence matrix with a root node row removed) is a basis if and only if the arcs corresponding to its columns form a spanning tree.%
\item \href{https://www.math.purdue.edu/~yipn/421/2024-11-05-NetFlows.pdf#page=18}{\textbf{Dual Variable Calculation from Spanning Tree}}: Given a spanning tree, dual variables are computed iteratively starting from a root node using the equation $y_j - y_i = c_\textbackslash{}{ij\textbackslash{}}$ for arcs $(i,j)$ in the spanning tree. Dual slacks for arcs not in the spanning tree are then computed using $z_\textbackslash{}{ij\textbackslash{}} = y_j - y_i + c_\textbackslash{}{ij\textbackslash{}}$.%
\item \href{https://www.math.purdue.edu/~yipn/421/2024-11-07-NetSimplex.pdf#page=45}{\textbf{Parametric Self-Dual Network Simplex Method}}: This method interweaves primal and dual pivots to reduce a perturbation parameter $\textbackslash{}mu$ from $\textbackslash{}infty$ to 0, finding an optimal solution by simultaneously improving primal and dual feasibility.%
\item \href{https://www.math.purdue.edu/~yipn/421/2024-11-14-NetAppls.pdf#page=17}{\textbf{Bellman-Ford Algorithm}}: $v_i = \textbackslash{}min_\textbackslash{}{j\textbackslash{}} \textbackslash{}\textbackslash{}{c_\textbackslash{}{ij\textbackslash{}} + v_j\textbackslash{}\textbackslash{}}$, $v_r = 0$%
\item \href{https://www.math.purdue.edu/~yipn/421/2024-11-19-MaxFlowMinCut.pdf#page=4}{\textbf{Max-Flow Min-Cut Theorem}}: The maximum flow from the source to the sink in a network is equal to the minimum capacity of any cut in the network.%
\end{itemize}

%
\subsection{Applications of Linear Programming}%
\label{subsec:ApplicationsofLinearProgramming}%
\begin{itemize}[leftmargin=*]%
\item \href{https://www.math.purdue.edu/~yipn/421/2024-10-28-ApplsLinProg.pdf#page=5}{\textbf{Ice Cream Production Planning}}: $\textbackslash{}min 50 \textbackslash{}sum_\textbackslash{}{i=1\textbackslash{}}^\textbackslash{}{12\textbackslash{}} (Y_i + Z_i) + 20 \textbackslash{}sum_\textbackslash{}{i=1\textbackslash{}}^\textbackslash{}{12\textbackslash{}} S_i$ s.t. $X_i + S_\textbackslash{}{i-1\textbackslash{}} - S_i = d_i$, $X_i - X_\textbackslash{}{i-1\textbackslash{}} = Y_i - Z_i$, $X_0 = 0$, $S_0 = 0$, $S_\textbackslash{}{12\textbackslash{}} = 0$, $X_i, Y_i, Z_i, S_i \textbackslash{}ge 0$%
\item \href{https://www.math.purdue.edu/~yipn/421/2024-10-28-ApplsLinProg.pdf#page=8}{\textbf{Workforce Planning Reformulation}}: $\textbackslash{}min 25 \textbackslash{}sum_\textbackslash{}{i=1\textbackslash{}}^\textbackslash{}{12\textbackslash{}} Y_i + (17.5 \textbackslash{}times 12) X$ s.t. $a_i - X \textbackslash{}le Y_i$, $0 \textbackslash{}le Y_i$, $X \textbackslash{}ge 0$%
\end{itemize}

%
\end{document}